\section{Rust 실행 구조와 유한 자원 계획}
\label{sec:runtime}

\subsection{형식화된 중간 표현과 소유권}

컴파일러는 구문과 이름 해석으로 연산자 그래프를 만들고,
형식 판단 $\Gamma\vdash p:\tau\xrightarrow{\mathcal E}\sigma$를 각 연결에 부여한다.
형식화된 중간 표현은 채널 스키마, 시계 사상, 발생 규칙,
품질 정책과 데이터 계보를 보유한다. 하드웨어 연결은 유효 설정과 capability 조건을,
실행 계획은 버퍼ㆍ호출시간ㆍ자원 배분 조건을 갖는다.
\begin{equation}
 \begin{aligned}
 \mathrm{Source}&\longrightarrow\mathrm{AST}
 \longrightarrow\mathrm{TypedGraph},\\
 \mathrm{TypedGraph}&\longrightarrow
        \mathrm{BoundGraph}(H,\theta)
 \longrightarrow\mathrm{ExecutionPlan}.
 \end{aligned}
 \label{eq:compiler-pipeline}
\end{equation}
각 산출물은 앞 단계의 노드 식별자와 계약을 유지한다.
격자ㆍ단위ㆍ시계 변환은 명시적 노드로 존재하며,
상태 초기화와 모델 활성화 사건은 실행 의존성에 포함된다.
코드 생성은 이 계획의 actor, 버퍼와 전이 관계를 Rust 구성 요소에 대응시킨다.

장치 actor는 native handle, SDK callback 등록, 설정과 획득 세대를 소유한다.
공통 수명 상태는 다음 전이로 구성한다.
\begin{equation}
 \begin{aligned}
 \mathsf{Created}&\longrightarrow\mathsf{Prepared}(\theta)
 \longrightarrow\mathsf{Streaming}(e),\\
 \mathsf{Streaming}(e)&\longrightarrow\mathsf{Draining}(e)
 \longrightarrow\mathsf{Closed}.
 \end{aligned}
 \label{eq:device-lifecycle}
\end{equation}
각 전이는 SDK의 완료 응답과 자원 소유 조건을 사용한다.
종료 전이는 callback 정리와 진행 중 호출의 완료를 포함한다.
설정 변경은 새 계약과 세대를 선택하는 제어 사건으로 작동한다.
fault 사건은 원인과 마지막 확정 상태를 보관하고,
복구 전이는 선택한 초기화 계약으로 새 세대를 생성한다.

가변 actor 상태는 한 소유자가 갱신하고,
완성된 데이터 블록은 소유권 이동 또는 읽기 공유로 전달한다.
\code{Arc<T>}를 사용하는 공유 객체는 $T$의 스레드 전달ㆍ공유 조건을 충족한다
\citep{rust_arc}. FFI wrapper는 포인터 유효성, 정렬, 길이, 수명과 aliasing 조건을
검증하여 native 버퍼를 형식화된 블록으로 연결한다\citep{rust_reference}.
이 경계별 검증 조건은 RustBelt가 제시한 의미론적 라이브러리 확장의
검증 구조를 따른다\citep{rustbelt}.

\subsection{고정 발생률 부분그래프의 실행 증명서}

한 설정과 스키마 세대에서 생산ㆍ소비 token 수가 일정한 부분그래프를
동기 데이터흐름 그래프 $G=(V,E)$로 나타낸다\citep{sdf}.
간선 $e$의 생산 수를 $p_e>0$, 소비 수를 $c_e>0$,
초기 token 수를 $d_e\ge0$으로 둔다.
양의 정수 반복벡터 $\boldsymbol\nu$는 균형식을 만족한다.
\begin{equation}
 p_e\nu_{\mathrm{src}(e)}
 =
 c_e\nu_{\mathrm{dst}(e)},
 \qquad e\in E.
 \label{eq:sdf-balance}
\end{equation}
순차 발화열 $\omega=(v_1,\ldots,v_M)$에서 $N_v(j)$는
첫 $j$회 발화 중 actor $v$의 횟수이다.
한 발화는 입력을 소비하고 상태 전이를 완료한 뒤 출력을 생산한다.
접두 발화열 이후의 점유량은
\begin{equation}
 m_e(j)=d_e+p_eN_{\mathrm{src}(e)}(j)
                 -c_eN_{\mathrm{dst}(e)}(j)
 \label{eq:sdf-occupancy}
\end{equation}
이다. 실행 증명서는 다음 조건을 모두 포함한다.
\begin{equation}
 \begin{aligned}
 N_v(M)&=\nu_v,\qquad v\in V,\\
 m_e(j-1)&\ge c_e,
       \qquad e:\mathrm{dst}(e)=v_j,\\
 B_e&\ge\max_{0\le j\le M}m_e(j).
 \end{aligned}
 \label{eq:sdf-certificate}
\end{equation}
입력 token 조건은 각 발화의 실행 가능성을 제공하고,
균형식은 $m_e(M)=d_e$를 준다.
각 actor의 전이가 유한 시간에 완료되고 외부 입력이 지정한 가용 조건을 만족하면,
실행기는 같은 발화열을 반복하며 용량 $B_e$ 안에서 흐름을 유지한다.
병렬 계획은 발화 시작 시 입력을 차감하고 완료 시 출력을 더하는 사건열을 사용하며,
허용한 사건열의 접두사마다 같은 token 조건과 용량 조건을 검사한다.

간선별 token 크기 상한이 $s_e$ byte이고 독립 저장 영역을 예약하면,
전체 메모리 예약량은 다음 식으로 계산한다.
\begin{equation}
 M_{\mathrm{reserve}}
 =
 \sum_{e\in E}B_es_e+
 \sum_{v\in V}M_v^{\mathrm{work}}+
 M^{\mathrm{meta}}.
 \label{eq:memory-reservation}
\end{equation}
$M_v^{\mathrm{work}}$는 actor가 동시에 소유하는 입력ㆍ출력ㆍ작업 영역의 상한이며,
$M^{\mathrm{meta}}$는 큐 관리와 제어 사건의 저장량이다.
공유 블록을 사용하는 계획은 공유 풀의 수명과 점유량으로 해당 항을 계산한다.
이벤트 marker와 설정 제어는 사건 actor로 처리하고,
고정률 영역의 경계에서 가용 조건과 세대 전이를 연결한다.

\subsection{호출시간과 주기 실행 조건}

순차 계획에 전용 실행 구간을 배정하고 actor $v$의 호출시간 상한을 $C_v$로 둔다.
발화 $j$의 의존 자료 가용 시각 상한을 $\overline a_j^{\mathrm{dep}}$,
호출ㆍ스케줄링 부가시간 상한을 $o_j$로 둔다.
주기 $k$의 시작이 $kP$일 때 완료 시각의 상한을 재귀적으로 계산한다.
\begin{equation}
 \begin{aligned}
 \overline f_{k,0}&=kP,\\
 \overline f_{k,j}
 &=\max\{\overline f_{k,j-1},
                 \overline a_{k,j}^{\mathrm{dep}}\}
                 +C_{v_j}+o_j.
 \end{aligned}
 \label{eq:cyclic-response}
\end{equation}
각 발화의 상대 deadline을 $D_j$로 두면,
\begin{equation}
 \overline f_{k,j}\le kP+D_j,
 \qquad
 \overline f_{k,M}\le(k+1)P
 \label{eq:cyclic-deadline}
\end{equation}
가 해당 순차 계획의 충분조건이다.
주기 시작 시 모든 외부 의존 자료가 준비되는 경우에는
$\sum_v\nu_vC_v+\sum_jo_j\le P$로 주기 조건을 계산한다.
실행 프로필은 $C_v$, $o_j$와 가용 상한의 근거,
CPU 배분과 OS 조건을 함께 보관한다.

\subsection{도착ㆍ서비스 곡선의 버퍼와 지연 상한}

전송 또는 처리 경계의 누적 유입량을 $A(t)$, 누적 완료량을 $D(t)$로 둔다.
두 양은 같은 payload 계량과 기준 시계를 사용한다.
초기 상태가 빈 FIFO에서 도착 곡선 $\alpha$와 서비스 곡선 $\beta$를 다음과 같이 정의한다.
\begin{equation}
 \begin{aligned}
 A(t)-A(s)&\le\alpha(t-s),\qquad 0\le s\le t,\\
 D(t)&\ge(A\otimes\beta)(t),\\
 (A\otimes\beta)(t)&=
       \inf_{0\le s\le t}\{A(s)+\beta(t-s)\}.
 \end{aligned}
 \label{eq:network-curves}
\end{equation}
burst 상한 $b$, 지속 유입률 상한 $r$,
해당 흐름에 배정한 유효 서비스율 $R>r$와 지연 $T$를 사용한다.
\begin{equation}
 \begin{aligned}
 \alpha(0)&=0,\qquad \alpha(u)=b+ru\quad(u>0),\\
 \beta(u)&=R(u-T)^+.
 \end{aligned}
 \label{eq:rate-latency}
\end{equation}
$r$는 기준 시계에서의 유입 상한이며, $b$는 패킷화와 burst를 포함한다.
$R,T$는 선택한 완료 경계까지의 자원 배분과 경합을 반영한다.
완성 블록을 계수하는 경계에서는 $T$에 블록 완성과 packetization의 지연을 포함한다.
Network calculus의 backlogㆍdelay 정리에 따라\citep{networkcalculus}
\begin{equation}
 \begin{aligned}
 A(t)-D(t)&\le b+rT,\\
 W(t)&=\inf\{d\ge0:A(t)\le D(t+d)\}
          \le T+\frac bR
 \end{aligned}
 \label{eq:network-bounds}
\end{equation}
를 얻는다. 동일한 payload 계량을 보존하는 직렬 경로는
$\beta_{\mathrm{path}}=\beta_1\otimes\cdots\otimes\beta_m$으로 합성하며,
rate--latency 모형에서는
$R_{\mathrm{path}}=\min_jR_j$,
$T_{\mathrm{path}}=\sum_jT_j$를 사용한다.
실행 계획은 이 상한을 ingress 이후의 가용 지연과 연결하여
식~\eqref{eq:window-availability-bound} 및 marker의 반응시간 조건을 계산한다.
